\subsection{Qualitative trajectory forecasts}
\label{sec:qualitative-trajectory-forecasts}

% Select one of the available instances by setting each flag to 01--05.
% Double pendulum uses two independent instances for its two regimes.
% These defaults may also be overridden before this file is input.
\providecommand{\DoublePendulumInstance}{01}
\providecommand{\DoublePendulumChaoticInstance}{02}
\providecommand{\DeterministicRosslerInstance}{01}
\providecommand{\DeterministicLorenzInstance}{01}
\providecommand{\StochasticRosslerInstance}{01}
\providecommand{\StochasticLorenzInstance}{01}

% Select five of the ten exported stochastic-Lorenz draws (01--10).
% Each selector is applied consistently to GT and every forecasting method.
\providecommand{\StochasticLorenzDrawOne}{01}
\providecommand{\StochasticLorenzDrawTwo}{02}
\providecommand{\StochasticLorenzDrawThree}{03}
\providecommand{\StochasticLorenzDrawFour}{04}
\providecommand{\StochasticLorenzDrawFive}{05}

We compare sampled forecast trajectories and Monte Carlo predictive means
across the evaluated systems and baselines.

\newlength{\qualpanelwidth}
\newlength{\qualpanelheight}
\newlength{\qualmethodswidth}
\setlength{\qualpanelwidth}{0.225\linewidth}
\setlength{\qualpanelheight}{0.215\linewidth}
\setlength{\qualmethodswidth}{0.735\linewidth}
\newcommand{\qualplaceholder}{%
  \fbox{\parbox[c][\dimexpr\qualpanelheight-2\fboxsep-2\fboxrule\relax][c]{%
    \dimexpr\qualpanelwidth-2\fboxsep-2\fboxrule\relax}{%
    \centering\scriptsize Image\\placeholder}}%
}
\newcommand{\qualgraphic}[1]{%
  \includegraphics[
    width=\qualpanelwidth,
    height=\qualpanelheight,
    keepaspectratio
  ]{#1}%
}
\newcommand{\qualpanel}[1]{%
  \IfFileExists{#1.pdf}{%
    \qualgraphic{#1.pdf}%
  }{%
    \IfFileExists{#1.png}{\qualgraphic{#1.png}}{\qualplaceholder}%
  }%
}
\newcommand{\qualheader}[1]{\scriptsize\bfseries\shortstack{#1}}
\newcommand{\qualcell}[2]{%
  \begin{minipage}[t]{\qualpanelwidth}
    \centering
    \qualheader{#1}\par\vspace{0.15em}
    \qualpanel{#2}
  \end{minipage}%
}
\newcommand{\qualcomparison}[2]{%
  \par\smallskip
  \noindent\textit{#1}\par\nobreak\vspace{0.15em}
  \noindent
  \begin{minipage}[c]{\qualpanelwidth}
    \centering
    \qualheader{GT}\par\vspace{0.15em}
    \qualpanel{#2_gt}
  \end{minipage}\hfill
  \begin{minipage}[c]{\qualmethodswidth}
    \centering
    \setlength{\tabcolsep}{1pt}
    \begin{tabular}{@{}ccc@{}}
      \qualcell{Block\\diffusion}{#2_block_diffusion}
      & \qualcell{AR\\diffusion}{#2_ar_diffusion}
      & \qualcell{DeepAR}{#2_deepar} \\[0.35em]
      \qualcell{Streaming\\flow}{#2_streaming_flow}
      & \qualcell{Chronos-2}{#2_chronos2}
      & \qualcell{VSF}{#2_vsf}
    \end{tabular}
  \end{minipage}\par
}
\newcommand{\qualinstancecomparison}[4]{%
  \qualcomparison{#1}{figures/#2/instance_#3_#4}%
}

\newcommand{\stochasticlorenzrow}[1]{%
  \scriptsize Draw~#1
  & \qualpanel{figures/lorenz_stochastic/instance_\StochasticLorenzInstance%
      _stochastic_lorenz_draw_#1_gt}
  & \qualpanel{figures/lorenz_stochastic/instance_\StochasticLorenzInstance%
      _stochastic_lorenz_draw_#1_block_diffusion}
  & \qualpanel{figures/lorenz_stochastic/instance_\StochasticLorenzInstance%
      _stochastic_lorenz_draw_#1_ar_diffusion}
  & \qualpanel{figures/lorenz_stochastic/instance_\StochasticLorenzInstance%
      _stochastic_lorenz_draw_#1_deepar}
  & \qualpanel{figures/lorenz_stochastic/instance_\StochasticLorenzInstance%
      _stochastic_lorenz_draw_#1_streaming_flow}
  & \qualpanel{figures/lorenz_stochastic/instance_\StochasticLorenzInstance%
      _stochastic_lorenz_draw_#1_vsf} \\
}

\begin{figure}[p]
  \centering
  \qualinstancecomparison{Sampled trajectories ($N=10$)}{double_pendulum}
    {\DoublePendulumInstance}{deterministic_double_pendulum_samples}
  \vspace{0.5em}
  \qualinstancecomparison{Predictive mean ($N=100$)}{double_pendulum}
    {\DoublePendulumInstance}{deterministic_double_pendulum_mean}
  \caption{Qualitative deterministic forecasts for the selected
  Double-pendulum instance.  GT is shown once on the left; the six forecasting
  methods are arranged in two rows of three.}
  \label{fig:qualitative-deterministic}
\end{figure}

\begin{figure}[p]
  \centering
  \qualinstancecomparison{Sampled trajectories ($N=10$)}{double_pendulum}
    {\DoublePendulumChaoticInstance}{deterministic_double_pendulum_samples}
  \vspace{0.5em}
  \qualinstancecomparison{Predictive mean ($N=100$)}{double_pendulum}
    {\DoublePendulumChaoticInstance}{deterministic_double_pendulum_mean}
  \caption{Qualitative deterministic forecasts for a more chaotic
  Double-pendulum instance.}
  \label{fig:qualitative-deterministic-double-pendulum-chaotic}
\end{figure}

\begin{figure}[p]
  \centering
  \qualinstancecomparison{Sampled trajectories ($N=10$)}{rossler_deterministic}
    {\DeterministicRosslerInstance}{deterministic_rossler_samples}
  \vspace{0.5em}
  \qualinstancecomparison{Predictive mean ($N=100$)}{rossler_deterministic}
    {\DeterministicRosslerInstance}{deterministic_rossler_mean}
  \caption{Qualitative deterministic forecasts for the selected R\"ossler
  instance.}
  \label{fig:qualitative-deterministic-rossler}
\end{figure}

\begin{figure}[p]
  \centering
  \qualinstancecomparison{Sampled trajectories ($N=10$)}{lorenz_deterministic}
    {\DeterministicLorenzInstance}{deterministic_lorenz_samples}
  \vspace{0.5em}
  \qualinstancecomparison{Predictive mean ($N=100$)}{lorenz_deterministic}
    {\DeterministicLorenzInstance}{deterministic_lorenz_mean}
  \caption{Qualitative deterministic forecasts for the selected Lorenz
  instance.}
  \label{fig:qualitative-deterministic-lorenz}
\end{figure}

\begin{figure}[p]
  \centering
  \qualinstancecomparison{Overlaid trajectories ($N=10$)}{rossler_stochastic}
    {\StochasticRosslerInstance}{stochastic_rossler_overlaid_draws}
  \caption{Conditional forecast distributions for the selected stochastic
  R\"ossler instance.  Each panel overlays ten trajectories sampled under the
  same observed context and forecast horizon.}
  \label{fig:qualitative-stochastic}
\end{figure}

\begin{figure}[p]
  \centering
  \begingroup
  \setlength{\qualpanelwidth}{0.136\linewidth}
  \setlength{\qualpanelheight}{0.13\linewidth}
  \setlength{\tabcolsep}{1.25pt}
  \renewcommand{\arraystretch}{1.05}
  \begin{tabular}{@{}l*{6}{c}@{}}
    & \qualheader{GT}
    & \qualheader{Block\\diffusion}
    & \qualheader{AR\\diffusion}
    & \qualheader{DeepAR}
    & \qualheader{Streaming\\flow}
    & \qualheader{VSF} \\
    \multicolumn{7}{@{}l}{\textbf{Stochastic Lorenz}} \\[-0.15em]
    \stochasticlorenzrow{\StochasticLorenzDrawOne}
    \stochasticlorenzrow{\StochasticLorenzDrawTwo}
    \stochasticlorenzrow{\StochasticLorenzDrawThree}
    \stochasticlorenzrow{\StochasticLorenzDrawFour}
    \stochasticlorenzrow{\StochasticLorenzDrawFive}
  \end{tabular}
  \endgroup
  \caption{Five selected draws from the conditional reference and forecast
  distributions for the selected stochastic Lorenz instance.  Each row uses
  the configurable draw index shown at left.}
  \label{fig:qualitative-stochastic-lorenz}
\end{figure}
